\scanpage{24}

\begin{center}$-\ 10\ -$\end{center}\hfill n\textsuperscript{o} 276

en vertu du corollaire précédent on peut donc se ramener au cas où $E$ est \uline{libre}, i.e. est produit direct d'un nombre fini de facteurs isomorphes à $A$. Raisonnant par récurrence, on est alors ramené à établir que si deux modules $E$ et $F$ (munis des filtrations associées à la filtration de $A$) sont complets il en est de même du module $E\times F$ ; or les relations
\[
A_pE\times A_qF\subset A_n(E\times F)\subset A_nE\times A_nF\qquad(p,q\leqslant n)
\]
montrent que $E\times F$ est le produit direct des groupes \uline{topologiques} $E$ et $F$. il est donc complet si $E$ et $F$ le sont.

PROPOSITION 5 -- \uline{Soient} $A$ \uline{un anneau commutatif}, $\mathrm m$ \uline{un idéal de} $A$, \uline{et supposons l'anneau} $A$ \uline{séparé} (resp. complet) \uline{pour la topologie} $\mathrm m$-\uline{adique}. \uline{Alors l'anneau de séries formelles} $A[[X]]$, \uline{filtré par les puissances de l'idéal} \inserted{\unsure{engendré}} \uline{par} $\mathrm m$ \uline{et} $X$, \uline{est séparé} (resp. \uline{complet}).
\margin{\unsure{Dire que la} \unsure{top. est} \unsure{celle de} la \unsure{conv. simple} des coefficients.}

Soit $\mathrm n$ l'idéal de $B=\unsure{'}A[[X]]$ engendré par $\mathrm m$ et $X$ : c'est l'ensemble des séries formelles $\sum a_nX^n$ telles que $a_0\in\mathrm m$. On déduit de là que pour tout entier $p$, l'idéal $\mathrm n^p$ est formé des séries $\sum a_nX^n$ telles que l'on ait $a_n\in\mathrm m^{p-n}$ pour $0\leqslant n\leqslant p$. Il résulte immédiatement de là que $B$ est séparé si $A$ est séparé. Supposons de plus $A$ complet, et considérons dans $B$ une série $\sum y_i$ dont le terme général
\[
y_i=\sum a_{in}X^n
\]
tend vers $0$ ; pour tout entier $p$ on a donc
\[
(1)\qquad a_{in}\in\mathrm m^{p-n}\qquad(0\leqslant n\leqslant p)\text{ pour }i>i_p.
\]
En fixant $n$ et en faisant varier $p$ on voit, puisque $A$ est complet, que la série $\sum_i a_{in}$ converge dans $A$ ; soit $a_n$ l'une des ses sommes (on ne suppose pas $A$ séparé). En vertu de (1) et du fait que les puissances de $\mathrm m$ sont fermées dans $A$, on a, quels que soient $n$ et $p\geqslant n$, la relation
\[
\sum_{i>i_p}\unsure{a_{in}|e_{in}}\in\mathrm m^{p-n}
\]
\mspage{23}
